\subsection{D1: localization, the density gap, and its exact repair}

\begin{statement}{D.1}{Proved}{The exact finite-descent localization map}
Let $C/K$ be smooth, proper and geometrically connected, of genus at
least two, over a number field. There is a bijection
\[
 \operatorname{loc}:\Sec(C/K)\xrightarrow{\ \sim\ }
 C(\mathbb A_K)_\bullet^{\mathrm{f\text{-}cov}}.
 \tag{D.1.1}
\]
Here the right side uses all finite-group-scheme torsors, twists by
global $H^1$, and connected components at the real places. If $K$
contains no CM subfield, its image in each individual $C(K_v)$,
for $v$ finite, is finite. The construction gives no uniform bound
as $v$ varies.
\end{statement}
\begin{proof}
The stipulated local theorem gives a unique point of $C(K_v)$
from the restriction of a global section at each finite place.
The real section theorem gives a unique real connected component;
at a complex place the Galois group is trivial. We use the
modified adelic set
\[
 C(\mathbb A_K)_\bullet=
 \prod_{v\nmid\infty}C(K_v)\times
 \prod_{v\text{ real}}\pi_0(C(K_v)).
\]
Properness makes every resulting tuple adelic. The hypotheses of
the local theorem and the real input are stated in
\cite[Theorem 1.1 and \S8]{AgentOpenAILocal}.

For a finite $K$-group scheme $G$ and a $G$-torsor $T\to C$,
survival means that the tuple lifts to $T^\xi$ for one
$\xi\in H^1(K,G)$. Harari--Stix's Theorem 2.1, with their
subgroup $U=1$ and omitted set $S=\varnothing$, identifies survival
of all such torsors with existence of a global section having the
prescribed local point sections. Their assertion (iii) uses
conjugacy by the geometric fundamental group. Their Remark 2.2(1)
identifies this condition with finite-cover descent
\cite[preprint Theorem 2.1, (i)$\Leftrightarrow$(iii), and
Remark 2.2(1)]{AgentHarariStix}. Torsor evaluations are constant
on real connected components, and complex places impose no
condition. The theorem therefore gives precisely the surjectivity
in (D.1.1), rather than a different quotient of the local data.

For injectivity, two classes with the same tuple have conjugate
restrictions at every finite place. Porowski's Theorem 1.1 on
locally conjugate Galois sections applies to a smooth curve over
a number field and to a density-one set of finite places. Take
that set to be all finite places. It gives geometric conjugacy of
the global sections \cite[Theorem 1.1]{AgentPorowskiConjugacy}.
This supplies the asserted injectivity independently of a
positive-density promotion of reduction equality.

Every section just considered is Selmer: its local restrictions
are point sections at all places, with the real-component
convention. Betts--Stix's Theorem A gives finiteness of the
individual localization image when $K$ has no CM subfield
\cite[Definition 1.1 and Theorem A]{AgentBettsStix}. The theorem
does not state a place-independent cardinality bound.
\end{proof}

\paragraph{Attribution of the gap.}
Section 8 of the supplied local-theorem paper was read in full.
Its Proposition 8.1 proves the image equality in (D.1.1), and its
last paragraph explicitly limits the Betts--Stix conclusion to
finite individual images. It contains no step replacing strongly
positive density by positive density
\cite[\S8, Proposition 8.1 and final paragraph]{AgentOpenAILocal}.
The text actually identified by LM6-D.6 is the unnumbered
corollary following Corollary 7.10 in Porowski's
\emph{Families of curves separating points}. The complete 2026
primary text was unavailable in the retrievals for this return;
the exact attribution is the one recorded in the supplied
LM6-D.6, rather than a new direct check of that primary corollary
\cite[D.5--D.6]{AgentLM6}. The prime-set assertion used there has
the following complete repair.

\begin{statement}{D.2}{Proved}{The strong-density criterion and the failed intersection}
Put $\mathcal S=\operatorname{Sel}(C/K)$ and
$r(v)=\#\operatorname{red}_v\mathcal S$ at the good finite places.
For $d\geq1$,
\[
 \#\mathcal S\leq d
 \quad\Longleftrightarrow\quad
 \text{there exists a strongly positive set $\Omega$ with
 $r(v)\leq d$ for all $v\in\Omega$}.
 \tag{D.2.1}
\]
Strongly positive means positive upper density after pullback to
every finite extension. A positive-density set can have empty
intersection with a potentially-density-one set, even over $\Q$.
\end{statement}
\begin{proof}
We use the exact reduction-rigidity statement supplied as
LM6-D.5: equality of the reductions of two Selmer sections on a
strongly positive set implies equality of their classes. Its
argument constructs a potentially-density-one separation set for
each fixed local partition, and uses the strong positivity to
meet that set \cite[D.5]{AgentLM6}. No version with an arbitrary
positive-density set is used here.

The forward implication of (D.2.1) holds with all good finite
places. For the reverse, take $d+1$ distinct classes. At every
$v\in\Omega$ some pair of them has equal reduction. Let
$T_1,\ldots,T_m$ be the finitely many corresponding pair-equality
sets; their union contains $\Omega$.

If each $T_i$ failed to be strongly positive, choose a finite
extension $L_i/K$ where its pullback has upper density zero, and
let $L$ be their compositum. Pullback to a further finite
extension preserves upper density zero. Here is the precise
counting reason. Primes of residue degree one over a prime $w$
have the same norm as $w$ and multiplicity at most the extension
degree. Primes of residue degree at least two and norm at most
$X$ lie over primes of norm at most $X^{1/2}$; their number is
$O(X^{1/2})$, negligible compared with the number of prime
ideals of norm at most $X$. The same estimates prove the claim
for upper Dirichlet density by summing $N(w)^{-s}$.
Thus every pullback $T_{i,L}$ has upper density zero. Their
finite union cannot contain the positive-upper-density set
$\Omega_L$. At least one $T_i$ is strongly positive, so reduction
rigidity identifies its pair, a contradiction.

For the failed intersection take
\[
 L=\Q(i),\qquad
 \Omega=\{p:p\equiv3\pmod4\},\qquad
 T=\{p:p\equiv1\pmod4\}.
 \tag{D.2.2}
\]
Both rational prime sets have density $1/2$. The pullback of
$T$ to $L$ has density one: the primes above $\Omega$ have norm
$p^2$, so their counting function is at most
$\pi(X^{1/2})$, whereas the finitely many ramified primes have
zero density. Thus $T$ is potentially density one, but
$\Omega\cap T=\varnothing$. This disproves exactly the
intersection inference. It does not construct a curve with more
than $d$ Selmer classes satisfying the ordinary-positive-density
bound.
\end{proof}

\paragraph{What depends on this step.}
The localization bijection D.1, local finiteness, and the
strongly positive statement of LM6-D.5 use no intersection with
an arbitrary positive-density set. LM6-D.6 itself identifies
the failed inference, and LM6-D.7 constructs an infinite compact
space with precisely the finite individual images and
potential-density partition refinements under discussion. The
local, finite-support and density-one integrality arguments in
AP-D.19, D.22 and D.25--D.26 likewise do not use that inference
\cite[D.5--D.7]{AgentLM6}\cite[D.19, D.22, D.25--D.26]{AgentAP}.
The result affected is the upgrade of the uniform criterion
from strongly positive to arbitrary positive density. Its
failure of proof is not a proof that its arithmetic conclusion
is false. In particular D1, finiteness of every proper
hyperbolic section set over a number field, remains undecided
in this return.

\subsection{Arithmetic counterexamples to the density promotion}

\begin{statement}{D.3}{Proved}{Actual Jacobians with equal positive-density traces and different potential toric ranks}
The genus-two curves
\[
 C_4:y^2=x^6+8x^4+16x^2+8,
 \qquad C_0:y^2=x^6+8
 \tag{D.3.1}
\]
have Jacobians of potential toric ranks $2$ and $0$ at $5$.
Their Frobenius traces agree at every good prime inert in
$\Q(\sqrt2)$, a density-$1/2$ set. Hence positive-density trace
equality does not force potential isogeny even for weight-one
Jacobian representations distinguished by potential toric rank.
There are also elliptic curves with full Frobenius-polynomial
equality on a density-$1/4$ set and no potential isogeny.
\end{statement}
\begin{proof}
For $a\in\Q\setminus\{-1,3\}$ set
\[
 \begin{split}
 E_a &: v^2=u^3+a u^2+a u+1,\\
 E_a^{(2)} &: 2v^2=u^3+a u^2+a u+1,\\
 C_a &: y^2=x^6+2a x^4+4a x^2+8.
 \end{split}
\]
The discriminant of the sextic is
$-2^{21}(a-3)^6(a+1)^2$, so the smooth proper completion has
genus two. The explicit maps
\[
 \pi_1:C_a\longrightarrow E_a^{(2)},\quad
 (x,y)\longmapsto(x^2/2,y/4),
 \qquad
 \pi_2:C_a\longrightarrow E_a,\quad
 (x,y)\longmapsto(2/x^2,y/x^3)
 \tag{D.3.2}
\]
satisfy the target equations by substitution. They extend to
the smooth proper curves by properness and have degree two.
Their invariant differentials pull back to
\[
 \pi_1^*(du/v)=4x\,dx/y,
 \qquad
 \pi_2^*(du/v)=-4\,dx/y.
\]
These are a basis of $H^0(C_a,\Omega^1)$. Consequently the
homomorphism
\[
 \Psi_a=(\pi_{1*},\pi_{2*}):J(C_a)\longrightarrow
 E_a^{(2)}\times E_a
 \tag{D.3.3}
\]
has invertible differential. Its connected image is a closed
subgroup of dimension two and its kernel has dimension zero;
thus it is a $\Q$-isogeny. This is the actual morphism relating
the Jacobian and the two elliptic factors. A quasi-inverse to
an isogeny has composites equal to multiplication by a nonzero
integer. After tensoring Tate modules with $\Q_\ell$, it
therefore gives an isomorphism of Galois representations.

Over $\Q(\sqrt2)$ the map $(u,v)\mapsto(u,\sqrt2v)$ identifies
$E_a^{(2)}$ with $E_a$. The nontrivial Galois element changes
this map by $[-1]$. Hence the induced Tate-module isomorphism
is
\[
 V_\ell(E_a^{(2)})\simeq V_\ell(E_a)\otimes\chi_2,
 \qquad
 a_p(J(C_a))=(1+\chi_2(p))a_p(E_a).
 \tag{D.3.4}
\]
Good-reduction Frobenius on the Tate module is the one whose
polynomial is obtained by point counting
\cite[V, Propositions 8.1 and 8.3]{AgentMilneEC}.
For $(2/p)=-1$ the last trace is zero independently of $a$.
The primes $p\equiv3,5\pmod8$ have density $1/2$.

For the displayed equation of $E_a$, the Weierstrass invariants
are
\[
 c_4=16a(a-3),\qquad
 \Delta=16(a-3)^3(a+1).
 \tag{D.3.5}
\]
At $5$, $E_4$ has $c_4=64$ and $\Delta=80$. Since $c_4$ is a
unit this equation is minimal there and its reduction is
multiplicative. Thus its potential toric rank is one.
For $E_0$, $\Delta=-432$ is a $5$-adic unit, giving good
reduction and potential toric rank zero. These reduction tests
are \cite[II, \S3]{AgentMilneEC}. The displayed twist
isomorphism makes the potential ranks of the twists equal to
the ranks of their original curves.

For completeness, the isogeny comparison preserves these
ranks: after a finite extension giving semistable reduction,
the logarithm $N$ of unipotent inertia on $V_\ell$, for
$\ell\ne5$, has rank equal to the toric rank. An isogeny
conjugates the rational Tate-module representations and their
operators $N$. Applying (D.3.3) gives
\[
 t_5(J(C_4))=1+1=2,\qquad t_5(J(C_0))=0+0=0.
 \tag{D.3.6}
\]
These ranks also persist after further finite extensions, since
the monodromy logarithm is multiplied by a nonzero ramification
index. No finite extension can therefore make the Jacobians
isogenous. Formula (D.3.4) proves the claimed trace equality.

We record precisely the invariant being compared. At an inert
good prime, writing $b=a_p(E_a)$ gives
\[
 P_{J(C_a),p}(T)
  =(T^2-bT+p)(T^2+bT+p)
  =T^4+(2p-b^2)T^2+p^2.
 \tag{D.3.7}
\]
At $p=19$ direct counting gives $a_{19}(E_4)=-4$ and
$a_{19}(E_0)=8$. The polynomials are respectively
$T^4+22T^2+361$ and $T^4-26T^2+361$.
Thus this pair proves the trace assertion; it does not prove
equality of their full characteristic polynomials on the inert
set.

For full-polynomial equality use
\[
 E_i:y^2=x^3-x,
 \qquad E_3:y^2=x^3+1.
 \tag{D.3.8}
\]
At $p\equiv11\pmod{12}$ both traces are zero. For $E_i$, pair
$x$ and $-x$ in the quadratic-character sum: $x^3-x$ changes
sign and $-1$ is a nonsquare. For $E_3$, cubing permutes
$\mathbf F_p$ when $p\equiv2\pmod3$, so its character sum is
the sum of the quadratic character over the entire field,
which is zero. Both characteristic polynomials are therefore
$T^2+p$. The prime set has density $1/4$.

These elliptic curves cannot become isogenous. At $5$ their
traces are $-2$ and $0$, with eigenvalue pairs
\[
 \{-1+2i,-1-2i\},\qquad
 \{i\sqrt5,-i\sqrt5\}.
\]
The first ratio is $(-3-4i)/5$, which is not a root of unity:
it belongs to $\Q(i)$ and differs from its four roots of unity.
The second ratio is $-1$. At a place of residue degree $f$
over $5$ in a finite extension, good-reduction Frobenius has
the $f$th powers of these eigenvalues. Equality of the two
unordered pairs would force the $f$th power of the first
ratio to be $(-1)^f$, or its inverse; that would make the first
ratio a root of unity. An isogeny would identify the rational
Tate modules and force this impossible equality. Finally an
isogeny over $\bar\Q$ is defined over a finite extension,
because its rational functions have finitely many algebraic
coefficients. This proves the last assertion.
\end{proof}

The all-prime identities in this proof follow from the displayed
isogenies and character sums. The program
\path{code/density_jacobian_check.py} checks the symbolic
substitutions, discriminant and Weierstrass invariants, the
two counts at $19$, and representative prime traces; its output
is retained in \path{code/density_jacobian_check_output.txt}.
These arithmetic certificates exclude promotion of the
potential-density trace argument to an ordinary-density-one
separation assertion for arbitrary such Jacobians. They are
not counterexamples to the Selmer-cardinality corollary in
D.2 or to finiteness in D1.

The hypotheses of available positive-density comparison theorems
matter here. Kulkarni--Patankar--Rajan retain a non-CM hypothesis
in the elliptic conclusion and a connected-monodromy hypothesis
in the relevant trace statement
\cite[Theorems 1.1, 1.5 and 1.6]{AgentKPR}.
Patankar--Rajan's corrected structural theorem retains its
stated additional monodromy hypotheses
\cite[Introduction and Example 1.1]{AgentPatankarRajan}.
The examples above give explicit maps and Frobenius calculations
showing why those hypotheses cannot simply be omitted in this
application. They do not preclude a different proof of D1.

\subsection{D2--D3: the closed-fibre test and generic reconstruction}

\begin{statement}{D.4}{Proved}{A Tate tower split at all closed points with no generic section}
For a finite extension $k/\Q_p$ there is a smooth hyperbolic
$k$-curve $B$ and a compatible tower of nontrivial connected
finite \'etale cyclic covers of $B$ split at every closed point.
Each nontrivial cover has no $k(B)$-point above the generic
point. More precisely, there is a continuous epimorphism
\[
 \theta:\pi_1(B)\longrightarrow\widehat\Z
 \tag{D.4.1}
\]
which is surjective on the geometric group and trivial on every
closed-point decomposition group. The tower does not construct
a nonpoint section of a non-isotrivial proper hyperbolic curve.
\end{statement}
\begin{proof}
Choose $q\in k$ with $0<|q|<1$, let $E_q$ be the Tate curve,
and put $B=E_q\setminus\{O\}$. For every finite extension
$L/k$ its uniformization is
\[
 E_q(L)=L^\times/q^{\Z}.
 \tag{D.4.2}
\]
For $n\ge1$ the identity on $L^\times$ induces the algebraic
isogeny
\[
 f_n:E_{q^n}\longrightarrow E_q,
 \qquad [a]_{q^{n\Z}}\longmapsto[a]_{q^\Z}.
 \tag{D.4.3}
\]
Its kernel is the constant cyclic group generated by $[q]$,
of order $n$, and its degree is $n$. These statements follow
from Tate's uniformization and isogeny theorem: in his notation
$t_2^m=t_1^n$ and the map is induced by $w\mapsto w^n$;
take $t_1=q^n,t_2=q,m=n,n=1$
\cite[Theorem 1 and the isogeny theorem on printed p.\ 16]{AgentTateUniformization}.
Set
\[
 B_n=E_{q^n}\setminus f_n^{-1}(O).
\]
Then $B_n\to B$ is connected, geometrically connected,
finite \'etale and cyclic of degree $n$.

Let $b$ be any closed point of $B$, with residue field
$L=k(b)$. Choose a representative $a\in L^\times$ for $b$
in (D.4.2). Its full fibre is
\[
 \{[a q^j]_{q^{n\Z}}:0\le j<n\}\subset B_n(L).
 \tag{D.4.4}
\]
All $n$ points are distinct, and none was removed because
$b\ne O$. Thus every closed fibre is split over its own
residue field.

For $n>1$, the extension of function fields
$k(B_n)/k(B)$ has degree $n$. A $k(B)$-point of its generic
fibre would give a $k(B)$-algebra homomorphism from that
field into $k(B)$. Such a homomorphism is injective and
$k(B)$-linear, contradicting the degree. This proves that
no open subscheme of $B$ admits a section of the cover.

For $n\mid m$, the identity on the uniformizing multiplicative
group gives compatible maps $B_m\to B_n$, carrying $[q]$
to $[q]$ on the deck groups. Choose $b_0\in B(k)$ and a
representative $a_0\in k^\times$. The points $[a_0]$ in all
levels choose compatible fibre basepoints. The resulting
quotient maps
$\theta_n:\pi_1(B)\to\Z/n\Z$ are compatible and yield
(D.4.1). Their geometric restrictions are surjective because
$B_{n,\bar k}$ is connected. Compactness shows that the
inverse-limit map on the geometric group is surjective: its
closed image surjects onto each finite quotient of
$\widehat\Z$.

By (D.4.4), the image of every closed-point decomposition
group in each cyclic deck group is trivial. Hence
$\theta$ vanishes on all such groups. This is an explicit
nonzero element of the kernel of
\[
 \operatorname{Hom}_{\mathrm{cts}}(\pi_1(B),\widehat\Z)
 \longrightarrow
 \prod_{b\in B_{\mathrm{cl}}}
 \operatorname{Hom}_{\mathrm{cts}}(G_{k(b)},\widehat\Z).
 \tag{D.4.5}
\]
Since the target is abelian, conjugacy choices for the
decomposition groups do not change this calculation.
\end{proof}

\paragraph{The reconstruction map that remains missing.}
For D3, the supplied LM6-B.38--B.40 give, using the local theorem,
the exact cartesian comparison for a smooth proper hyperbolic
family $\mathcal X\to B$ over a smooth $p$-adic base:
\[
 \pi_1(X_{k(B)})\simeq
 \pi_1(\mathcal X)\times_{\pi_1(B)}G_{k(B)}.
 \tag{D.4.6}
\]
The generic section has first component killing the
codimension-one inertia defining $\pi_1(B)$. It consequently
gives a splitting
\[
 \sigma:\pi_1(B)\longrightarrow\pi_1(\mathcal X)
 \quad\text{of}\quad
 \pi_1(\mathcal X)\longrightarrow\pi_1(B).
 \tag{D.4.7}
\]
Restriction at a closed point $b$ is the section of a unique
point of $\mathcal X_b(k(b))$ by the local theorem. These are
the exact maps proved in R12; they do not supply a morphism
$B\to\mathcal X$. In the isotrivial case the target comparison
becomes a map to a fixed hyperbolic curve, and LM6-B.41 applies
its geometric-image trichotomy and curve-target Hom theorem
to produce that morphism \cite[B.38--B.42]{AgentLM6}.
For a non-isotrivial family the constant target used by that
argument has not been constructed. D.4 verifies explicitly
that closed-fibre splitting alone cannot replace the missing
generic argument over a $p$-adic base. Its relative dimension
is zero, so it is not presented as a counterexample to D3's
proper hyperbolic relative curve.

For D2, Bresciani calls the corresponding generic reconstruction
property \emph{coherence}. The published version proves
coherence for a constant abelian family by first applying
Faltings's Tate-module Hom theorem and then using finite
generation of Mordell--Weil groups to recover an integral
Hom class \cite[Proposition 4.5]{AgentBrescianiFiniteType}.
His Lemma 4.7 propagates it through a finite morphism using
Hilbert irreducibility over a Hilbertian field. Corollary 4.8
then treats a family with a finite morphism into a constant
abelian scheme. These are the precise morphisms behind the
constant-abelian-cover cases in Theorem B and the
hyperelliptic and genus-at-most-two consequences. The
unrestricted coherence proof in Proposition 4.3 invokes the
additional weak Bombieri--Lang input in Theorem 4.2; the paper
does not eliminate that input. D.4 applies over a $p$-adic
field and does not contradict the Hilbertian hypothesis of
Lemma 4.7. Neither the unrestricted reconstruction in D2
nor the non-isotrivial assertion in D3 is decided here.

\subsection*{Audit of the stipulated local theorem}

Theorem 1.1 of \cite{AgentOpenAILocal} is an unrefereed
preprint input explicitly authorized by the prompt. We checked
the internal steps used here and the exact hypotheses of the
principal imported geometric statements. No specific error was
found in those checked steps. This audit does not certify every
external theorem or every argument of that preprint. In
particular the resolution-of-nonsingularities input is the
authors' revised RIMS preprint, not a statement whose
publication status has been inferred from its use here.

\paragraph{The analytic point and all inverse branches.}
In Sections 2--3 the completions of the algebraic fields embed
isometrically into the fixed completion of the separable
closure. The resulting $H$-valued analytic point retains its
residue-field embedding. One cannot replace it merely by its
underlying point over the smaller constant field. Near a
horizontal branch point the local equation can be put in the
form $t=z^e$ after shrinking. For a unit $1+h$, the formula
\[
 (1+h)^{1/e}=\exp\bigl(e^{-1}\log(1+h)\bigr)
\]
converges whenever the unit is sufficiently close to $1$;
the bound $v_p(h)>v_p(e)+1/(p-1)$ suffices. This includes
$p\mid e$. Properness and compactness let one shrink the
base disk so that its entire inverse image is the union of
the specified branch disks.

The patching proposition requires a root on \emph{every}
inverse collar. When $m$ is divisible by each $e$, these
roots split every factor of $Z^e-t$ on that collar. Distinct
roots differ by a nonzero constant times an invertible
analytic function, so the factors are comaximal even when
$p\mid e$. One glues the pulled-back finite cover on the
exterior to trivial finite covers on the inner disks. The
finite analytic algebra then algebraizes by proper analytic
GAGA. A nonempty connected component of the resulting finite
\'etale cover surjects onto the connected base. The remaining
analytic map on the exterior is evaluated at the given
$H$-point; the proof does not require that exterior map to
descend to an algebraic map over the smaller field.

\paragraph{Stable edges, the lifting theorem, and the full collar.}
The Section 4 circuit construction uses a prime $\ell\ne p$.
At a circuit vertex choose the cyclic residue-curve cover
ramified exactly at the two circuit marks; divisibility of
the degree-zero Picard group by $\ell$ gives the required
rational function with divisor $P-Q$ modulo $\ell$. All
other marks are unramified and split over the algebraically
closed residue field. Edge expansion is $\ell$ on the
circuit and one on its other incident edges, with the
required copies on the latter. This satisfies the local
degree sums and the tame Riemann--Hurwitz condition. A
separating edge is first made nonseparating by connected
unramified cyclic covers on its two stable sides, joining
their lifted attachment points. Each side has a cycle or
a positive-genus vertex: a tree of $N$ rational vertices
with only its one attaching leg has total valence $2N-1$,
incompatible with stability. Thus the required nontrivial
unramified cyclic cover exists on each side. The successive
fibre products preserve the root on every inverse branch.

The exact imported lifting statement is Amini--Baker--Brugall\'e--Rabinoff
Theorem 7.7 for a tame cover of metrized complexes over a fixed
triangulated punctured curve
\cite[Definition 4.31 and Theorem 7.7]{AgentABBR}.
Their definition places the branch locus inside the selected
puncture set. Taking that set empty therefore yields an
\'etale curve cover, as required here.

For the reduction from arbitrary collars, the precise
resolution input is Mochizuki--Tsujimura's Theorem A for a
hyperbolic curve over a finite $p$-adic field with a prime
set $\Sigma$ containing $p$ and at least one other prime,
and their Proposition 2.4(iv)
\cite[Theorem A and Proposition 2.4(iv)]{AgentMochizukiTsujimura}.
The preprint takes $\Sigma$ to be all primes. After finite
constant extension, the stable model of a connected
geometrically pro-$\Sigma$ finite \'etale Galois cover maps
to the chosen compactified model. The asserted model
morphism need not be finite; the argument uses its
compatibility with reduction. On an open-disk tube an
invertible analytic pullback coordinate has constant
absolute value in the value group of the completed
algebraic constants. An irrational logarithmic radius
therefore avoids all such tubes and the model vertices.
Its full inverse level lies in node annuli. Compactness
then places a smaller entire inverse collar in finitely
many of these annuli, where the stable-edge lemma applies.

\paragraph{Arbitrary roots and horizontal ramification.}
Section 5 labels the relevant connected components of a
torus-uniformized inverse image by
$\Lambda/\Lambda^m$. Adjusted torus lifts glue because the
lattice $\Lambda$ has no $m$-torsion. This is used to get
roots over every inverse component. In the auxiliary
cyclic cover the dual graph has two rational vertices and
$\ell$ parallel edges; a cycle $e_j-e_{j'}$ gives slope
one on the selected collar. The middle-collar Laurent
coefficient bounds put the remaining analytic unit in the
logarithm/exponential range above. Horizontal ramification
is measured in characteristic-zero closed-point local
fields $\bar k((t))$. Their ramification indices in a
compositum divide the relevant least common multiple;
normalization after a base change of index $r$ changes an
index $e$ to $e/\gcd(e,r)$. An \'etale cover on top adds
no horizontal index. The proof uses precisely these
characteristic-zero facts, rather than a claim of tame
ramification at residue characteristic $p$.

\paragraph{The field embedding and final decomposition group.}
In Section 6 a base-field embedding of a finite Galois
subextension $M$ into $H$, inside the fixed ambient
completion, has image $M$: all roots of each minimal
polynomial already belong to $M$. The embedding constructed
there therefore proves the required subfield containment.
In Section 7 restriction of the decomposition groups is
surjective by transitivity of prolongations and injective
by the preceding completion result. Its inverse is
continuous by compactness. The argument excluding two
independent henselian valuations uses
\[
 q_r=d+(1-d)\frac{z^r}{1+z^r}.
\]
For sufficiently large $r$, one valuation makes $q_r/d$
a square and the other makes $q_r$ a square. Their quotient
would make the chosen nonsquare constant $d$ a square.
The estimates in that argument include $p=2$.

The remaining rank-two valuation has horizontal
coarsening $\operatorname{ord}_a$ and residue field $k(a)$.
Surjectivity of the section onto $G_k$ forces $k(a)=k$.
The horizontal inertia maps trivially to the \'etale
fundamental group of the original proper curve, so the
section obtained is its point section. For injectivity,
the Jacobian Kummer argument uses all multiplication maps,
including the $p$-power ones on the compact $p$-adic Lie
group. Dropping the $p$-power factors would not prove that
step. These checks justify the specific uses of the local
theorem in D.1 and (D.4.6)--(D.4.7) at the stated
preprint level; they do not settle the global finiteness
or generic-reconstruction questions.
